\section{Adversarial Pressure, Bias y Non-Authority}

Los usuarios reales no se ciñen a la definición de la tarea como lo haría un benchmark annotator. En cuanto se publica una demo, los usuarios buscan los límites, plantean comparaciones entre dos males, sustituyen términos de identidad y combinan a propósito situaciones poco comunes. Esta sección trata estos ataques como evidence, no como un episodio de PR.

\subsection{Exposed Socket: también hay que explicar por qué un ejemplo es correcto}

La página de noticias describe el reto de la moneda en el enchufe que Alexa le propuso a una niña; a la derecha, Delphi juzga que «dejar que un niño toque un enchufe expuesto es peligroso». Este ejemplo muestra que el sistema puede relacionar el lenguaje natural con consecuencias físicas, pero es solo un caso de éxito. Una evaluation real necesita reformular el agente de la acción, el tiempo, la negación y el hypothetical context, y comprobar si el modelo conserva el danger concept.

\lecturefigure{slide-36-alexa-socket-delphi.jpg}{El incidente del enchufe expuesto de Alexa y el juicio de peligro de Delphi}{Diapositiva V036 recuperada del video oficial; Stanford CS25 Lecture 16}

\subsection{Horas después del lanzamiento, los usuarios reescriben la definición de la tarea}

La backlash slide registra los cambios tras la publicación del 15 de octubre: el relative comparison mode se retiró en cuestión de horas, en un fin de semana se recibieron unos 25,000 adversarial examples y después las consultas reales llegaron a varios millones. Al leer la figura, estas cifras deben entenderse como attack-surface evidence; no equivalen a un conjunto de prueba uniforme, pero revelan un espacio de combinaciones que un benchmark difícilmente puede enumerar de antemano.

\lecturefigure{slide-37-delphi-backlash-scale.jpg}{Retiro rápido de Delphi tras su lanzamiento, 25k adversarial examples y tráfico real}{Diapositiva V037 recuperada del video oficial; Stanford CS25 Lecture 16}

La siguiente página muestra cómo investigadores y público construyen a propósito oraciones anómalas para «torturar» al modelo. Este tipo de ejemplos suele carecer de contexto realista, pero es muy eficaz para exponer lexical shortcuts. Al evaluar no hay que elegir entre «la pregunta es absurda, así que el fallo no cuenta» y «un solo fallo vuelve inútil todo el sistema»; lo adecuado es construir una failure taxonomy: negation, role reversal, rare composition, identity substitution y nonsensical premise.

\lecturefigure{slide-38-adversarial-twitter-example.jpg}{Contrived adversarial prompt en redes sociales}{Diapositiva V038 recuperada del video oficial; Stanford CS25 Lecture 16}

\subsection{La Public Critique no es ruido que se pueda ignorar}

El artículo de \emph{The New York Times} «Can a Machine Learn Morality?», con su ilustración de un dragón, llevó la controversia al ámbito público. La página siguiente enumera las concerned voices y vuelve al título del artículo académico. Juntas, ambas páginas muestran que quienes evalúan la machine ethics no son solo los ML reviewers, sino también especialistas en ética, periodistas, los grupos afectados por el modelo y los usuarios comunes; la documentación del sistema debe permitir que todas estas personas vean el objetivo de entrenamiento y sus limitaciones.

\lecturefigure{slide-39-can-machine-learn-morality-nyt.jpg}{La pregunta de los medios: ¿puede una máquina aprender moral?}{Diapositiva V039 recuperada del video oficial; Stanford CS25 Lecture 16}

\lecturefigure{slide-40-concerned-voices.jpg}{Preocupaciones públicas sobre Delphi y reflexión de los investigadores}{Diapositiva V040 recuperada del video oficial; Stanford CS25 Lecture 16}

\subsection{Delphi no es una Moral Authority}

La non-authority slide lo dice explícitamente: el equipo de investigación no respalda el uso de IA para dar moral advice, y Delphi no sustituirá a los human judges en los tribunales; la cita de abajo explica que el objetivo del proyecto es entender cómo los imperfect humans emiten juicios éticos, no convertir a la máquina en la autoridad moral última. Este límite debe incorporarse a toda licencia de uso posterior y a la UI, no quedarse solo en la discussion del artículo.

\lecturefigure{slide-41-no-ai-moral-authority.jpg}{Límite explícito: Delphi no da moral advice ni sustituye a los jueces humanos}{Diapositiva V041 recuperada del video oficial; Stanford CS25 Lecture 16}

\begin{importantbox}{La Capability Claim y la Authority Claim deben separarse}
Que un modelo pueda predecir una parte de los juicios humanos es una capability claim; que el modelo deba decidir con base en ello castigos, tratamientos médicos, contrataciones o resultados judiciales es una authority claim. Aunque la primera sea cierta, de ella no se deduce automáticamente la segunda. La autoridad requiere un diseño de instituciones, responsabilidad, apelación y auditoría.
\end{importantbox}

\subsection{«No investigar» también tiene una Normative Consequence}

La página siguiente reconoce primero que los moral models son demasiado difíciles y que su precisión es limitada, y luego señala que los sistemas de IA actuales ya toman decisiones éticas implícitas. El callout final propone invertir en investigación y formar normas sociales human-centered, unbiased y robust. Al leer la figura hay que considerar los riesgos de ambos lados: desplegar un imperfect moral model causa daño, pero no investigar en absoluto también puede dejar en los sistemas reglas predeterminadas que nadie ha auditado.

\lecturefigure{slide-42-ethics-investment-needed.jpg}{Que sea difícil no significa que se pueda ignorar: los sistemas actuales ya toman decisiones con consecuencias éticas}{Diapositiva V042 recuperada del video oficial; Stanford CS25 Lecture 16}

«Delphi will empower powerful people and harm marginalized people» es la preocupación de la postura contraria; el build posterior señala que los sistemas existentes ya pueden reinforce unjust status quo. Lo importante no es cuál de las burbujas gana, sino que la evaluación del despliegue debe comparar la intervention con la baseline: ¿cuál es el daño incremental que causa el nuevo sistema y a quién está dañando el proceso actual?

\lecturefigure{slide-43-reinforce-status-quo.jpg}{Comparación en ambos sentidos entre intervention risk y status-quo risk}{Diapositiva V043 recuperada del video oficial; Stanford CS25 Lecture 16}

\subsection{La Bias Acknowledgement debe concretarse en Data y Use}

La bias slide indica explícitamente que el norm bank refleja las ideologies de la annotation data de 2020--2021, que los judgments los aportaron crowdworkers de Estados Unidos y se inclinan hacia Western-centric viewpoints, y que adaptar Delphi a otras culturas o moral frameworks es una línea abierta. El titular de la izquierda, «Delphi leans left», recuerda que la inclinación política también se mide y se difunde.

\lecturefigure{slide-44-delphi-bias-acknowledgement.jpg}{Declaración de sesgos de datos, políticos y culturales de Delphi}{Diapositiva V044 recuperada del video oficial; Stanford CS25 Lecture 16}

\teachervoice{La ponente no trata la bias acknowledgement como una declaración que exime de culpa. Afirma que Delphi efectivamente es left-leaning y Western-centric, y luego orienta la discusión hacia una participación más amplia y el intercambio de datos. Reconocer el sesgo es el punto de partida de la auditoría, no una licencia para desplegar.}

\subsection{Resumen de la sección}

La demo pública llevó a Delphi de un benchmark estático a un entorno adversarial, político e institucional. El límite más importante de la clase es la non-authority: predecir juicios descriptivos no otorga poder de decisión. Una evaluación confiable debe registrar a la vez los fallos del modelo, los ataques de los usuarios, los grupos de los que provienen los datos y la status-quo baseline.
